% GENERATED FILE. Edit canonical lessons, then run npm run generate:books.
% canonical-source-hash: fnv1a64:795df0265d5a71aa
% canonical-lessons: HI-C299-angara, HI-C299-sulagna, HI-C299-sola, HI-C299-bharakna, HI-C299-phus

\chapter{Fire!}
\label{ch:hi-fire}
\hlchaptermodality{299}

\begin{chapteropening}
\textbf{By the end of this chapter:} \emph{I can say where a fire started and how it spread.}

Complete the last lesson of chapter 299: I can say where a fire started and how it spread.
\end{chapteropening}

\section[\texorpdfstring{\hi{अंगारा} (aṅgārā)}{अंगारा (aṅgārā)}]{\hi{अंगारा} (aṅgārā) — an ember; a live coal}

\label{lesson:HI-C299-angara}



\begin{quote}
Before the new one: say the Hindi for to crush, then the Hindi for to steady oneself.
\end{quote}

\subsection*{You'll want to know: \hi{अंगारा}}
\textbf{\hi{अंगारा}} — \emph{aṅgārā} — \textquotedblleft{}an ember; a live coal\textquotedblright{}.

\begin{quote}
\hi{चूल्हे} \hi{से} \hi{एक} \hi{अंगारा} \hi{गिरा।}
\end{quote}

\begin{quote}
\emph{cūlhe se ek aṅgārā girā.}
\end{quote}

\begin{quote}
An ember fell out of the stove.
\end{quote}

Masculine.

\begin{cousinweb}
comes from Sanskrit \emph{aṅgāra}, charcoal, burning or not (Monier-Williams). One small \textbf{\hi{अंगारा}} is enough to start a fire.
\end{cousinweb}

\begin{grammarlens}[title={What you've built}]
One more word for telling what happened.
\end{grammarlens}

\subsection*{Your turn}
\begin{itemize}
\raggedright
  \item \emph{Say it:} \emph{aṅgārā}
  \item \emph{Say it:} \emph{aṅgārā}, once more
  \item \emph{Say it:} \emph{cūlhe se ek aṅgārā girā.} — An ember fell out of the stove.
  \item \emph{Recall:} say the Hindi for to crush, then the Hindi for to steady oneself, then say \emph{aṅgārā} again
\end{itemize}

\begin{tcolorbox}[breakable,colback=teal!4,colframe=teal!35!black,title={Before you move on}]
What does \hi{अंगारा} mean? (\textquotedblleft{}An ember; a live coal\textquotedblright{}.) Say it once more.
\end{tcolorbox}
\section[\texorpdfstring{\hi{सुलगना} (sulagnā)}{सुलगना (sulagnā)}]{\hi{सुलगना} (sulagnā) — to smoulder}

\label{lesson:HI-C299-sulagna}



\begin{quote}
Before the new one: say the Hindi for to steady oneself, then the Hindi for an ember.
\end{quote}

\subsection*{You'll want to know: \hi{सुलगना}}
\textbf{\hi{सुलगना}} — \emph{sulagnā} — \textquotedblleft{}to smoulder\textquotedblright{}.

\begin{quote}
\hi{अंगारे} \hi{से} \hi{लकड़ी} \hi{सुलग} \hi{रही} \hi{थी।}
\end{quote}

\begin{quote}
\emph{aṅgāre se lakṛī sulag rahī thī.}
\end{quote}

\begin{quote}
The wood was smouldering from the ember.
\end{quote}

Past (it …): \hi{वह} \hi{लकड़ी} \hi{सुलगी}, \emph{vah lakṛī sulgī}.

\begin{cousinweb}
\textbf{\hi{सुलगना}} (\emph{sulagnā}) is set by CDIAL under Sanskrit \emph{saṁlagna}, \textquotedblleft{}stuck together\textquotedblright{} (CDIAL 12999), where the Hindi verb means to catch fire. It is the slow burning that comes before the flames.
\end{cousinweb}

\begin{grammarlens}[title={What you've built}]
One more word for telling what happened.
\end{grammarlens}

\subsection*{Your turn}
\begin{itemize}
\raggedright
  \item \emph{Say it:} \emph{sulagnā}
  \item \emph{Say it:} \emph{sulagnā}, once more
  \item \emph{Say it:} \emph{aṅgāre se lakṛī sulag rahī thī.} — The wood was smouldering from the ember.
  \item \emph{Recall:} say the Hindi for to steady oneself, then the Hindi for an ember, then say \emph{sulagnā} again
\end{itemize}

\begin{tcolorbox}[breakable,colback=teal!4,colframe=teal!35!black,title={Before you move on}]
What does \hi{सुलगना} mean? (\textquotedblleft{}To smoulder\textquotedblright{}.) Say it once more.
\end{tcolorbox}
\section[\texorpdfstring{\hi{शोला} (śolā)}{शोला (śolā)}]{\hi{शोला} (śolā) — a flame; a burst of fire}

\label{lesson:HI-C299-sola}



\begin{quote}
Before the new one: say the Hindi for an ember, then the Hindi for to smoulder.
\end{quote}

\subsection*{You'll want to know: \hi{शोला}}
\textbf{\hi{शोला}} — \emph{śolā} — \textquotedblleft{}a flame; a burst of fire\textquotedblright{}.

\begin{quote}
\hi{खिड़की} \hi{से} \hi{आग} \hi{के} \hi{शोले} \hi{निकल} \hi{रहे} \hi{थे।}
\end{quote}

\begin{quote}
\emph{khiṛkī se āg ke śole nikal rahe the.}
\end{quote}

\begin{quote}
Flames were leaping out of the window.
\end{quote}

Masculine.

\begin{cousinweb}
\textbf{\hi{शोला}} (\emph{śolā}) is usually traced, through Persian, to Arabic \emph{shuʿla}, \textquotedblleft{}a flame, a blaze\textquotedblright{}. A burning house sends out \textbf{\hi{शोले}}; so, in speech, do angry eyes.
\end{cousinweb}

\begin{grammarlens}[title={What you've built}]
One more word for telling what happened.
\end{grammarlens}

\subsection*{Your turn}
\begin{itemize}
\raggedright
  \item \emph{Say it:} \emph{śolā}
  \item \emph{Say it:} \emph{śolā}, once more
  \item \emph{Say it:} \emph{khiṛkī se āg ke śole nikal rahe the.} — Flames were leaping out of the window.
  \item \emph{Recall:} say the Hindi for an ember, then the Hindi for to smoulder, then say \emph{śolā} again
\end{itemize}

\begin{tcolorbox}[breakable,colback=teal!4,colframe=teal!35!black,title={Before you move on}]
What does \hi{शोला} mean? (\textquotedblleft{}A flame; a burst of fire\textquotedblright{}.) Say it once more.
\end{tcolorbox}
\section[\texorpdfstring{\hi{भड़कना} (bhaṛaknā)}{भड़कना (bhaṛaknā)}]{\hi{भड़कना} (bhaṛaknā) — to flare up}

\label{lesson:HI-C299-bharakna}



\begin{quote}
Before the new one: say the Hindi for to smoulder, then the Hindi for a flame.
\end{quote}

\subsection*{You'll want to know: \hi{भड़कना}}
\textbf{\hi{भड़कना}} — \emph{bhaṛaknā} — \textquotedblleft{}to flare up\textquotedblright{}.

\begin{quote}
\hi{हवा} \hi{से} \hi{आग} \hi{और} \hi{भड़क} \hi{गई।}
\end{quote}

\begin{quote}
\emph{havā se āg aur bhaṛak gaī.}
\end{quote}

\begin{quote}
The wind made the fire flare up even more.
\end{quote}

Past (it …): \hi{वह} \hi{आग} \hi{भड़की}, \emph{vah āg bhaṛkī}.

\begin{cousinweb}
\textbf{\hi{भड़कना}} (\emph{bhaṛaknā}) is inherited from \emph{*bhaṭ}, a sudden jerk or crackle, a form CDIAL reconstructs (CDIAL 9365). A fire \textbf{\hi{भड़कती} \hi{है}}, and so can anger or a crowd.
\end{cousinweb}

\begin{grammarlens}[title={What you've built}]
One more word for telling what happened.
\end{grammarlens}

\subsection*{Your turn}
\begin{itemize}
\raggedright
  \item \emph{Say it:} \emph{bhaṛaknā}
  \item \emph{Say it:} \emph{bhaṛaknā}, once more
  \item \emph{Say it:} \emph{havā se āg aur bhaṛak gaī.} — The wind made the fire flare up even more.
  \item \emph{Recall:} say the Hindi for to smoulder, then the Hindi for a flame, then say \emph{bhaṛaknā} again
\end{itemize}

\begin{tcolorbox}[breakable,colback=teal!4,colframe=teal!35!black,title={Before you move on}]
What does \hi{भड़कना} mean? (\textquotedblleft{}To flare up\textquotedblright{}.) Say it once more.
\end{tcolorbox}
\section[\texorpdfstring{\hi{फूस} (phūs)}{फूस (phūs)}]{\hi{फूस} (phūs) — dry straw; thatch}

\label{lesson:HI-C299-phus}



\begin{quote}
Before the new one: say the Hindi for a flame, then the Hindi for to flare up.
\end{quote}

\subsection*{You'll want to know: \hi{फूस}}
\textbf{\hi{फूस}} — \emph{phūs} — \textquotedblleft{}dry straw; thatch\textquotedblright{}.

\begin{quote}
\hi{फूस} \hi{की} \hi{छत} \hi{पल} \hi{भर} \hi{में} \hi{जल} \hi{गई।}
\end{quote}

\begin{quote}
\emph{phūs kī chat pal bhar meṁ jal gaī.}
\end{quote}

\begin{quote}
The thatched roof burned in a moment.
\end{quote}

Masculine.

\begin{cousinweb}
\textbf{\hi{फूस}} (\emph{phūs}) is inherited from \emph{*phussa}, \textquotedblleft{}defective\textquotedblright{}, a form CDIAL reconstructs (CDIAL 9100), under which it gathers words for old dry straw, chaff and rubbish. \textbf{\hi{फूस} \hi{की} \hi{छत}} is a thatched roof.
\end{cousinweb}

\begin{grammarlens}[title={What you've built}]
One more word for telling what happened.
\end{grammarlens}

\subsection*{Your turn}
\begin{itemize}
\raggedright
  \item \emph{Say it:} \emph{phūs}
  \item \emph{Say it:} \emph{phūs}, once more
  \item \emph{Say it:} \emph{phūs kī chat pal bhar meṁ jal gaī.} — The thatched roof burned in a moment.
  \item \emph{Recall:} say the Hindi for a flame, then the Hindi for to flare up, then say \emph{phūs} again
\end{itemize}

\begin{tcolorbox}[breakable,colback=teal!4,colframe=teal!35!black,title={Before you move on}]
What does \hi{फूस} mean? (\textquotedblleft{}Dry straw; thatch\textquotedblright{}.) Say it once more.
\end{tcolorbox}
